\section{模型检验与评价}
\subsection{问题一的物理约束与费用复核}
使用未舍入调度序列计算各区间的供需平衡残差和储能递推残差：
\begin{align}
 r_k^{\mathrm{bal}}&=g_k+d_k+u_k-\ell_k-c_k,\\
 r_k^{\mathrm{soc}}&=E_{k+1}-E_k-0.9c_k+d_k/0.9.
\end{align}
同时核对购电与光伏消纳边界、充放电功率、互斥条件、全部145个状态及日初日末条件。电量和费用一致性分别采用$10^{-6}$ kWh、$10^{-6}$元的绝对容差，复核结果见表\ref{tab:validation}。
\begin{table}[!htbp]
\centering
\caption{问题一求解结果的独立复核}
\label{tab:validation}
\small
\begin{tabular}{lll}
\toprule
检验项目 & 计算结果 & 判定 \\
\midrule
最大供需平衡残差绝对值 & $1.14\times10^{-13}$ kWh & 通过 \\
最大储能递推残差绝对值 & $9.09\times10^{-13}$ kWh & 通过 \\
复算费用与求解目标之差 & 0元 & 通过 \\
储电量与充放电功率边界 & 均满足约束 & 通过 \\
逐区间充放电互斥 & 未出现同时充放电 & 通过 \\
购电非负与光伏消纳边界 & 均满足约束 & 通过 \\
日初、日末储电量 & 均为6\,000 kWh & 通过 \\
\bottomrule
\end{tabular}
\end{table}
供需与储能残差均低于设定容差，全天能量收支也满足式\eqref{eq:roundtrip}和式\eqref{eq:dailybalance}。求解器目标值与下界相等，则进一步给出了该模型下数值最优性的依据。

\subsection{问题四的连续回放检验}
问题四在两类合同条件下分别核对点预测基准与情景采购，各轨迹覆盖48\,096个区间，均从同一储电量独立运行。逐区间复核实际供需平衡、连续状态、动作保护、充放电互斥和合同权限，并由合同及交易记录重新结算费用。最新两条情景轨迹均通过完整回放检查，最大状态连接差为0，实际末态分别为$5999.999999999999$、$6000.0$ kWh，均满足$10^{-6}$ kWh绝对容差；最后一日全部145个状态满足6\,000 kWh预留下界。

除实际轨迹外，复核还覆盖每次滚动求解的完整三情景计划：独立检查各情景物理方程、合同权限、共同动作、加权目标和复用后缀误差界。以决策时刻为界重放可见历史及已发布预报，重建原点预测、已实现误差库和代表情景，并核对计划与执行记录。两采购方案每类条件的1\,336次预测发布中，原负载、光伏、价格及其来源逐项相同。因此，费用差对应情景规划机制的改变，不能据此报告点预测精度改善。

冷启动期两类情景方案各运行1\,008步，与点预测基准动作及费用一致；从初态连续两周各2\,016步，覆盖冷启动及情景启用阶段，通过物理及费用检查。另在2月8日12:00受控中断并恢复，预测、合同、执行、账单及每日汇总与对应连续运行一致，支持状态传递与误差库重建的正确性。两份年度Excel分别读回65\,782、114\,639个单元格，各覆盖48\,096个源时隙。

这些检查分别针对实际执行、情景计划和导出文件，并非新的参数扰动实验。当前年度已用于开发诊断与候选比较，决策阶段的时间因果正确性不能证明跨年份泛化能力，尚需独立年份评价。

\subsection{模型评价}
模型的主要优点在于电量定义明确：母线侧交换量和电池内部储电量分别进入供需平衡与状态递推，可直接核对转换损耗。单日模型通过初末状态相等条件排除初始储能净消耗的影响；连续模型则保留实际状态，区分购电合同、可用信息和实际账单，便于解释预测误差如何传导至调度费用。

模型也存在三方面限制。其一，常数效率和十分钟平均功率难以描述设备效率曲线及瞬时峰值。其二，费用中未计电池寿命损耗，情景采购增加储能吞吐，加入寿命成本后应重新评价。其三，三个代表情景及其固定权重只近似近期误差，且全窗口电池动作共享限制了未来调节能力；现有同条件回放不能据此宣称跨年份稳健性。最后一日预留提高了储能恢复余量，但不保证任意实际供需序列都可行。
\FloatBarrier
